\chapter{Phased System Roadmap}
\label{ch:roadmap}

The system is built in five phases. Phase~0 prepares the environment and
data. Phases~1 to~4 each deliver something usable on its own: a still-image
detector, a video analyser, a live camera system and anonymous multi-camera
analytics. Each phase ends with an exit test. The next phase starts only when
the test passes. Figure~\ref{fig:phases} summarises the sequence.

\begin{figure}[htbp]
\centering
\begin{tikzpicture}[
    pc/.style={base, text width=2.05cm, minimum height=1.32cm, font=\scriptsize},
    rl/.style={grp, font=\scriptsize\scshape, text width=1.55cm, align=right, anchor=east},
    ph/.style={grp, font=\footnotesize\scshape, text width=2.3cm},
  ]
  \def\dx{2.48}
  \def\x#1{{#1*\dx}}
  \def\ra{0}\def\rb{-1.52}\def\rc{-3.04}\def\rd{-4.56}\def\re{-6.08}

  % phase headings and arrows between them
  \foreach \i/\lab in {0/{Phase 0\\environment}, 1/{Phase 1\\still images}, 2/{Phase 2\\uploaded video}, 3/{Phase 3\\live streams}, 4/{Phase 4\\multi-camera}}
    \node[ph] (h\i) at (\i*\dx,1.2) {\lab};
  \foreach \i/\j in {0/1,1/2,2/3,3/4}
    \draw[arr] ($(h\i.east)+(-0.05,0)$) -- ($(h\j.west)+(0.05,0)$);

  % row labels
  \node[rl] at (-1.2,\ra) {input};
  \node[rl] at (-1.2,\rb) {core work};
  \node[rl] at (-1.2,\rc) {tools};
  \node[rl] at (-1.2,\rd) {deliverable};
  \node[rl] at (-1.2,\re) {exit test};

  % phase 0
  \node[pc, n] at (0*\dx,\ra) {laptop, Kaggle,\\ dataset requests};
  \node[pc, n] at (0*\dx,\rb) {WSL2, pinned\\ environment,\\ scan check};
  \node[pc, n] at (0*\dx,\rc) {uv or conda,\\ pytest, Kaggle\\ datasets};
  \node[pc, n] at (0*\dx,\rd) {reproducible\\ environment};
  \node[pc, n] at (0*\dx,\re) {tests pass\\ in both places};

  % phase 1
  \node[pc, fuse] at (1*\dx,\ra) {folders of\\ paired images};
  \node[pc, fuse] at (1*\dx,\rb) {baselines,\\ CMFM blocks,\\ ablations};
  \node[pc, fuse] at (1*\dx,\rc) {Ultralytics 8.4,\\ COCO eval,\\ SAHI, FiftyOne};
  \node[pc, fuse] at (1*\dx,\rd) {still-image\\ \method, tables};
  \node[pc, fuse] at (1*\dx,\re) {fusion gap\\ reproduced;\\ three seeds};

  % phase 2
  \node[pc, tim] at (2*\dx,\ra) {video files,\\ paired clips};
  \node[pc, tim] at (2*\dx,\rb) {clip-mode TSM,\\ tracking,\\ writers};
  \node[pc, tim] at (2*\dx,\rc) {PyAV, trackers,\\ TrackEval,\\ Gradio};
  \node[pc, tim] at (2*\dx,\rd) {video analyser\\ and app};
  \node[pc, tim] at (2*\dx,\re) {matches offline;\\ HOTA end to end};

  % phase 3
  \node[pc, rgb] at (3*\dx,\ra) {webcam, RTSP,\\ thermal camera};
  \node[pc, rgb] at (3*\dx,\rb) {capture threads,\\ TensorRT,\\ stream mode};
  \node[pc, rgb] at (3*\dx,\rc) {OpenCV,\\ TensorRT 11,\\ FastRTC};
  \node[pc, rgb] at (3*\dx,\rd) {live detection\\ system};
  \node[pc, rgb] at (3*\dx,\re) {30\,FPS for\\ ten minutes;\\ recovers};

  % phase 4
  \node[pc, dat] at (4*\dx,\ra) {calibrated\\ camera network};
  \node[pc, dat] at (4*\dx,\rb) {ground plane,\\ association,\\ zones, events};
  \node[pc, dat] at (4*\dx,\rc) {supervision\\ zones,\\ DeepStream 9.1};
  \node[pc, dat] at (4*\dx,\rd) {anonymous\\ analytics};
  \node[pc, dat] at (4*\dx,\re) {counts verified;\\ privacy review\\ signed};
\end{tikzpicture}

\vspace{2mm}
\caption[The five phases of the build]{The five phases of the build. Each
column is a phase and each row one aspect of it. A phase begins only when
the exit test of the previous phase has passed. Phases 1 to 3 share one set
of weights. Phase 4 is optional and bounded by the privacy rules of
Chapter~\ref{ch:ethics}.}
\label{fig:phases}
\end{figure}


\section{Phase 0: environment and data access}
\label{sec:phase0}

Phase~0 removes the obstacles that stall projects later.

\paragraph{Local machine}
On the Windows laptop, WSL2 with Ubuntu 24.04 provides a Linux environment.
GPU access comes from the Windows NVIDIA driver; no Linux driver is
installed inside WSL. Datasets are kept on the Linux file system, not under
\texttt{/mnt/c}, because cross-system file access is slow. A single Python
environment pins PyTorch, Ultralytics 8.4, the evaluation tools and the
video libraries. Jupyter runs inside WSL and is opened from the Windows
browser or from VS Code.

\paragraph{Kaggle}
A setup notebook installs the same pinned environment and compiles
\texttt{mamba\_ssm} once. It saves the wheel as a private dataset and
verifies the kernels on the T4. It also checks that the pure-PyTorch scan
agrees with the CUDA scan.

\paragraph{Data requests}
Requests for RGBT-Tiny, LLVIP, FLIR ADAS and KAIST are sent on the first
day, because approval can take weeks.

\paragraph{Exit test}
Phase~0 is complete when four checks pass.
\begin{itemize}
  \item The unit tests pass both locally and on Kaggle.
  \item The two scan implementations agree to within $10^{-4}$.
  \item A small YOLO26 trains for one epoch on a VisDrone subset in both
        places.
  \item The first datasets are packaged.
\end{itemize}

\section{Phase 1: still images}
\label{sec:phase1}

Phase~1 delivers the still-image \method{}, the baselines it is compared
with, and a tool that runs any of them on a folder of images.
Figure~\ref{fig:phase1} shows the pipeline. Images are read in pairs and
letterboxed to the model's input. For large aerial images, they are
optionally sliced with SAHI~\citep{akyon2022sahi}. Predictions are written
as COCO JSON and scored externally, and the results are inspected
visually in FiftyOne.

The order of work is set by risk.
\begin{enumerate}
  \item Single-modality YOLO26 baselines are trained first. They test the
        whole pipeline on a known model.
  \item Simple two-stream fusion follows.
  \item The CMFM block is added last, once the pipeline around it is
        trusted.
\end{enumerate}

\paragraph{Exit test}
The two-stream baseline must reproduce the published gap between
single-modality and fused detection on LLVIP or M3FD within two points. The
still-image \method{} must then complete three seeds on RGBT-Tiny, with
results generated by Notebook~15.

\begin{figure}[htbp]
\centering
\begin{tikzpicture}[
    pn/.style={base, text width=2.05cm, minimum height=1.4cm, font=\scriptsize},
  ]
  \def\dx{2.62}
  \def\ya{0}\def\yb{-2.35}

  \node[pn, dat]  (a) at (0*\dx,\ya) {\textbf{image folder}\\ paired files,\\ equal names};
  \node[pn, n]    (b) at (1*\dx,\ya) {\textbf{paired loader}\\ letterbox to\\ $640{\times}512$};
  \node[op, fill=cNeutral!14, minimum size=9mm, font=\scriptsize, align=center] (q) at (2*\dx,\ya) {large\\ image?};
  \node[pn, n]    (c) at (3*\dx,{\ya+1.05}) {\textbf{SAHI slicing}\\ 256 to 640\,px tiles,\\ 20\% overlap};
  \node[pn, fuse] (d) at (4*\dx,\ya) {\textbf{\method}\\ image mode,\\ PyTorch or TRT};
  \node[pn, n]    (e) at (4*\dx,\yb) {\textbf{merge}\\ tile predictions,\\ NMS or NMM};
  \node[pn, n]    (f) at (3*\dx,\yb) {\textbf{predictions}\\ COCO JSON\\ per image};
  \node[pn, alert](g) at (2*\dx,\yb) {\textbf{scoring}\\ AP, AP\textsubscript{S},\\ AI-TOD bins, SAFit};
  \node[pn, fuse] (h) at (1*\dx,\yb) {\textbf{tables}\\ generated by\\ Notebook 15};
  \node[pn, tim]  (i) at (0*\dx,\yb) {\textbf{inspection}\\ FiftyOne, error\\ categories};

  \draw[arr] (a) -- (b);
  \draw[arr] (b) -- (q);
  \draw[arr] (q) |- node[lbl, pos=0.3, left=1pt] {yes} (c);
  \draw[arr] (c) -| (d);
  \draw[arr] (q) -- node[lbl, above=0.5pt] {no} (d);
  \draw[arr] (d) -- (e);
  \draw[arr] (e) -- (f);
  \draw[arr] (f) -- (g);
  \draw[arr] (g) -- (h);
  \draw[arr] (f.south) -- ++(0,-0.45) -| (i.south);
\end{tikzpicture}

\vspace{2mm}
\caption[Phase 1 pipeline for still images]{Phase 1 pipeline for still
images. Large images are optionally sliced, and tile predictions are merged
before scoring. Predictions are exported as COCO JSON and scored with
external tools, so that size-specific and SAFit metrics are available.
Errors are inspected and categorised in FiftyOne before any change to the
model.}
\label{fig:phase1}
\end{figure}


\section{Phase 2: uploaded video}
\label{sec:phase2}

Phase~2 turns the detector into a video analyser. A user uploads a video,
or a pair of visible and thermal videos. The system returns an annotated
video, a track file and summary statistics (Figure~\ref{fig:phase2}).

\paragraph{Decoding}
Videos are decoded with PyAV, which ships Windows wheels with FFmpeg
included, or with OpenCV~\citep{pyav_repo}. The \texttt{decord} library is
avoided, because its last release dates from 2021. Paired videos are matched
frame by frame using their timestamps. Each thermal frame is assigned to
the nearest visible frame within half a frame interval. Unmatched frames
are recorded as dropouts for the stream monitor, not silently discarded.

\paragraph{Clip-mode inference}
\methodS{} runs causally from the first frame and carries its state through
the video. The same video can be re-run in image mode. This gives a direct,
per-video comparison of what the memory contributes.

\paragraph{Tracking}
Detections are linked into tracks with the trackers built into Ultralytics
(ByteTrack, BoT-SORT and TrackTrack) or with the Apache-licensed Roboflow
\texttt{trackers} package~\citep{roboflowtrackers_repo}. Camera-motion
compensation in BoT-SORT matters for drone footage. Tracks are written in
MOTChallenge format and scored with TrackEval.

\paragraph{Interface}
A Gradio application accepts uploads and shows the annotated result beside
the per-frame latency~\citep{gradio_repo}.

\paragraph{Exit test}
On the RGBT-Tiny and VT-Tiny-MOT test sequences, the pipeline must produce
the same detections as the offline evaluation code. Its tracking metrics
must be computed end to end by TrackEval.

\begin{figure}[htbp]
\centering
\begin{tikzpicture}[
    pn/.style={base, text width=2.05cm, minimum height=1.4cm, font=\scriptsize},
  ]
  \def\dx{2.62}
  \def\ya{0}\def\yb{-2.35}

  \node[pn, dat]  (a) at (0*\dx,\ya) {\textbf{upload}\\ one video or a\\ visible--thermal pair};
  \node[pn, n]    (b) at (1*\dx,\ya) {\textbf{decoder}\\ PyAV or OpenCV,\\ timestamps kept};
  \node[pn, n]    (c) at (2*\dx,\ya) {\textbf{pairing}\\ nearest frame\\ within $\tfrac12$ interval};
  \node[pn, dat]  (d) at (3*\dx,\ya) {\textbf{stream monitor}\\ $\tau_t$, dropout\\ flags, cuts};
  \node[pn, tim]  (e) at (4*\dx,\ya) {\textbf{\methodS}\\ clip mode, state\\ carried forward};

  \node[pn, n]    (f) at (4*\dx,\yb) {\textbf{tracker}\\ BoT-SORT with\\ motion compensation};
  \node[pn, fuse] (g) at (3*\dx,\yb) {\textbf{writers}\\ annotated MP4,\\ MOT tracks, CSV};
  \node[pn, alert](h) at (2*\dx,\yb) {\textbf{evaluation}\\ TrackEval: HOTA,\\ IDF1, flicker};
  \node[pn, fuse] (i) at (1*\dx,\yb) {\textbf{Gradio app}\\ result beside\\ frame latency};
  \node[pn, n]    (j) at (0*\dx,\yb) {\textbf{comparison}\\ same video in\\ image mode};

  \draw[arr] (a) -- (b); \draw[arr] (b) -- (c); \draw[arr] (c) -- (d); \draw[arr] (d) -- (e);
  \draw[arr] (e) -- (f); \draw[arr] (f) -- (g); \draw[arr] (g) -- (h); \draw[arr] (h) -- (i);
  \draw[darr] (i) -- (j);
  \draw[darr, draw=cAlert] (c.south) -- node[lbl, right=1pt, pos=0.45, text=cAlert] {unmatched frames become dropouts} ++(0,-0.62);
\end{tikzpicture}

\vspace{2mm}
\caption[Phase 2 pipeline for uploaded video]{Phase 2 pipeline for uploaded
video. Decoding keeps timestamps so that visible and thermal frames can be
paired and the true interval measured. Unmatched frames are recorded as
dropouts, not discarded. The streaming model runs causally from the first
frame. The same video is re-run in image mode to show what the memory adds.}
\label{fig:phase2}
\end{figure}


\section{Phase 3: live camera streams}
\label{sec:phase3}

Phase~3 connects the streaming model to cameras: a webcam, an IP camera over
RTSP, and a USB or network thermal camera. The difficulty here is not
accuracy but time. Figure~\ref{fig:phase3} shows the threaded design.

\paragraph{Capture}
Each camera has its own capture thread. On Windows, webcams are opened
through the DirectShow or Media Foundation back ends. RTSP streams are read
over TCP to avoid packet loss. Each thread writes into a queue that holds a
single frame: the newest frame replaces the old one, so the model always
sees the present and never a backlog. Every frame carries its capture
timestamp. The stream monitor pairs visible and thermal frames by timestamp
and computes the true interval $\tau_t$. It raises the sensor flags when a
camera stalls.

\paragraph{Inference}
The inference worker holds \methodS{} as a TensorRT FP16 engine in stream
mode. Its state persists between calls. Motion estimation runs in a separate
CPU thread on downscaled frames, in parallel with the network. The tracker
and renderer run in a third thread. The annotated stream reaches the browser
by WebRTC through FastRTC, or by MJPEG on a local
network~\citep{fastrtc_repo}.

\paragraph{Budget}
At 30 frames per second each frame has about 33\,ms
(Table~\ref{tab:latency}). Because the stages run in parallel, throughput
is limited by the slowest stage, not by the sum. End-to-end delay, from
light entering the camera to pixels on the screen, is the sum. The target is
under 150\,ms. When the model cannot keep up, three strategies apply, in
this order.
\begin{enumerate}
  \item Frames are dropped at capture, never in the middle of the pipeline.
  \item The interval factor $\tau_t/\tau_0$ tells the TSM how much time has
        passed.
  \item The detector can run every second frame, with the tracker predicting
        in between.
\end{enumerate}

\begin{table}[htbp]
\centering
\caption[Latency budget for live operation]{Latency budget for live operation
at 30 frames per second. Values other than the YOLO26 reference are planning
estimates to be measured in Notebook~13.}
\label{tab:latency}
\small
\begin{tabularx}{\textwidth}{@{}LR{2.6cm}P{4.6cm}@{}}
\hd{Stage} & \hd{Estimate (ms)} & \hd{Notes} \\ \gap
capture and decode & 5--10 & per camera thread, in parallel \\
pre-processing & 1--3 & letterbox and normalise on the GPU \\
network, two backbones, CMFM and TSM & 8--15 & YOLO26-s alone is 2.5\,ms on a T4~\citep{yolo26_2026} \\
NMS & 0--1 & none with the one-to-one head \\
motion estimation & 2--4 & CPU thread, overlaps the network \\
tracking & 1--3 & BoT-SORT or ByteTrack \\
rendering and encoding & 5--10 & overlay and WebRTC encoding \\
network to browser & 20--60 & local network or internet \\ \gap
\textit{end to end} & \textit{under 150} & \textit{target} \\
\end{tabularx}
\end{table}

\paragraph{Exit test}
On the desktop GPU, the live system must sustain 30 frames per second at
$640\times512$ with paired cameras for ten minutes, without growth in delay
or memory. It must also recover correctly when a camera is unplugged and
reconnected. The same test on Jetson Orin at 15 frames per second answers
H5.

\begin{figure}[htbp]
\centering
\begin{tikzpicture}[
    pn/.style={base, text width=2.15cm, minimum height=1.35cm, font=\scriptsize},
    qn/.style={base, fill=cNeutral!7, text width=0.95cm, minimum height=0.75cm, font=\scriptsize},
  ]
  % thread lanes
  \node[grp, anchor=west] at (-1.2,1.95) {capture threads};
  \node[pn, rgb] (cv) at (0,0.75)  {\textbf{visible camera}\\ webcam or\\ RTSP over TCP};
  \node[pn, ir]  (ct) at (0,-1.15) {\textbf{thermal camera}\\ USB or IP\\ module};
  \node[qn] (qv) at (2.2,0.75)  {queue\\ of 1};
  \node[qn] (qt) at (2.2,-1.15) {queue\\ of 1};
  \draw[arr] (cv) -- (qv);
  \draw[arr] (ct) -- (qt);
  \node[lbl, anchor=north] at (1.1,-2.2) {newest frame replaces the old};

  \node[grp, anchor=west] at (3.35,1.95) {inference thread};
  \node[pn, dat] (mon) at (4.5,-0.2) {\textbf{stream monitor}\\ pairs by time,\\ $\tau_t$ and flags};
  \draw[arr] (qv.east) -- (mon.west |- qv.east);
  \draw[arr] (qt.east) -- (mon.west |- qt.east);

  \node[pn, tim, minimum height=1.6cm] (inf) at (7.15,-0.2) {\textbf{\methodS}\\ TensorRT FP16,\\ stream mode,\\ state persists};
  \draw[arr] (mon) -- (inf);

  \node[pn, n] (mot) at (7.15,-2.55) {\textbf{motion thread}\\ ORB, RANSAC\\ on small frames};
  \draw[arr] (mot) -- node[lbl, right=1pt] {$H_{t-1\to t}$} (inf);
  \draw[arr] (mon.south) |- (mot.west);

  \node[grp, anchor=west] at (8.75,1.95) {output thread};
  \node[pn, n]    (trk) at (9.85,-0.2) {\textbf{tracker and}\\ \textbf{renderer}\\ boxes and tracks};
  \draw[arr] (inf) -- (trk);
  \node[pn, fuse] (web) at (9.85,-2.55) {\textbf{WebRTC}\\ FastRTC, or\\ MJPEG on LAN};
  \draw[arr] (trk) -- (web);
  \node[pn, fuse, text width=1.45cm] (br) at (12.2,-2.55) {\textbf{browser}\\ dashboard};
  \draw[arr] (web) -- (br);

  % latency strip
  \def\k{0.118}
  \def\yl{-4.95}
  \def\xs{-1.2}
  \node[grp, anchor=west] at (\xs,{\yl+0.95}) {one frame, end to end, planning estimate in milliseconds};
  \foreach \a/\w/\c/\lab in {0/8/cRGB/capture, 8/2/cNeutral/pre, 10/12/cTime/network, 22/2/cNeutral/track, 24/8/cFuse/render, 32/40/cData/{to browser}}{
    \fill[\c!32] ({\xs+\a*\k},{\yl-0.2}) rectangle ({\xs+(\a+\w)*\k},{\yl+0.2});
  }
  \foreach \a/\lab in {4/capture, 16/network, 28/render, 52/{network to browser}}
    \node[font=\scriptsize, text=ink, anchor=north] at ({\xs+\a*\k},{\yl-0.25}) {\lab};
  \foreach \t in {0,10,22,32,72}{
    \draw[draw=cNeutral, line width=0.45pt] ({\xs+\t*\k},{\yl+0.26}) -- ({\xs+\t*\k},{\yl+0.38});
    \node[font=\scriptsize, text=muted, anchor=south] at ({\xs+\t*\k},{\yl+0.38}) {\t};
  }
  \node[lbl, anchor=west] at ({\xs+72*\k+0.2},\yl) {target under 150};
\end{tikzpicture}

\vspace{2mm}
\caption[Phase 3 threaded design for live streams]{Phase 3 threaded design
for live streams. Each camera has a capture thread writing into a queue of
one frame, so the model always processes the present. The stream monitor
pairs frames by timestamp and supplies $\tau_t$ and the sensor flags.
Motion estimation runs in its own thread, in parallel with the network.
The strip below shows a planning estimate of where the time goes for one
frame. Stages run concurrently, so throughput is set by the slowest stage
and delay by their sum.}
\label{fig:phase3}
\end{figure}


\section{Phase 4: anonymous multi-camera analytics}
\label{sec:phase4}

Phase~4 is optional. It is the closest the project comes to the fictional
Machine, and it is deliberately limited (Chapter~\ref{ch:ethics}). Several
calibrated cameras each run Phase~3. Their tracks are projected onto a common
ground plane through camera-to-ground homographies. Tracks are associated
across cameras by position, time and motion. For vehicles, non-biometric
appearance such as colour and shape can be added. Person tracks are never
associated through faces or other biometric features. The outputs are
aggregate and anonymous:
\begin{itemize}
  \item counts per zone;
  \item line crossings;
  \item dwell times;
  \item heat maps;
  \item rule-based events such as intrusion into a restricted area or a
        vehicle travelling the wrong way.
\end{itemize}
On NVIDIA hardware, DeepStream~9.1 offers multi-camera pipelines and
three-dimensional tracking that can replace the custom
orchestration~\citep{deepstream_docs}. Figure~\ref{fig:phase4} shows the
design.

\paragraph{Exit test}
Faces are blurred at ingest. No identity or face embedding is stored, and
recordings are deleted on schedule. The counts must agree with manual
counts on a held-out recording. The privacy review of
Section~\ref{sec:design-rules} must be signed off.

\begin{figure}[htbp]
\centering
\begin{tikzpicture}[
    pn/.style={base, text width=2.2cm, minimum height=1.35cm, font=\scriptsize},
    cam/.style={base, rgb, text width=1.85cm, minimum height=0.95cm, font=\scriptsize},
  ]
  % cameras, each running phase 3
  \foreach \i/\y in {1/1.35, 2/0, 3/-1.35}
    \node[cam] (c\i) at (0,\y) {camera \i\\ Phase 3 pipeline};
  \node[font=\scriptsize, text=muted] at (0,-2.2) {$\vdots$};

  \node[pn, alert] (blur) at (2.75,0) {\textbf{ingest guard}\\ faces blurred,\\ no identity kept};
  \foreach \i in {1,2,3} \draw[arr] (c\i.east) -- (blur.west);

  \node[pn, n]    (gp)  at (5.45,0) {\textbf{ground plane}\\ camera-to-ground\\ homographies};
  \node[pn, tim]  (as)  at (8.15,0) {\textbf{association}\\ position, time,\\ motion; vehicle\\ appearance only};
  \node[pn, fuse] (an)  at (10.85,0) {\textbf{analytics}\\ zones, lines,\\ dwell, heat maps};
  \draw[arr] (blur) -- (gp); \draw[arr] (gp) -- (as); \draw[arr] (as) -- (an);

  \node[pn, fuse] (ev)  at (10.85,-2.45) {\textbf{events}\\ intrusion,\\ wrong-way travel};
  \node[pn, n]    (st)  at (8.15,-2.45) {\textbf{storage}\\ aggregates only;\\ raw video expires};
  \node[pn, dat]  (db)  at (5.45,-2.45) {\textbf{dashboard}\\ counts and alerts\\ for operators};
  \draw[arr] (an) -- (ev);
  \draw[arr] (ev) -- (st);
  \draw[arr] (st) -- (db);

  \node[pn, alert, text width=2.6cm] (pol) at (8.15,2.15) {\textbf{privacy rules}\\ no watch-lists, no face\\ embeddings, retention timer};
  \draw[darr, draw=cAlert] (pol.west) -| (blur.north);
  \draw[darr, draw=cAlert] (pol.south) -- (as.north);
  \draw[darr, draw=cAlert] (pol.east) -- (12.55,2.15) -- (12.55,-3.6) -| (st.south);
\end{tikzpicture}

\vspace{2mm}
\caption[Phase 4 anonymous multi-camera analytics]{Phase 4 anonymous
multi-camera analytics. Every camera runs the Phase 3 pipeline. An ingest
guard blurs faces and discards identity before anything is stored. Tracks
are associated across cameras on a shared ground plane by position, time and
motion. Appearance cues are used only for vehicles. The outputs are
aggregates and rule-based events. The privacy rules (red, dashed) constrain
ingest, association and storage.}
\label{fig:phase4}
\end{figure}


\section{Exit criteria at a glance}

Table~\ref{tab:exit} collects the exit test and the deliverable of every
phase in one place.

\begin{table}[htbp]
\centering
\caption[Exit criteria for each phase]{Exit criteria for each phase.}
\label{tab:exit}
\small
\begin{tabularx}{\textwidth}{@{}P{2.2cm}LP{3.3cm}@{}}
\hd{Phase} & \hd{Exit test} & \hd{Deliverable} \\ \gap
0, environment & tests pass locally and on Kaggle; scans agree to $10^{-4}$; data packaged & reproducible environment \\
1, still images & published fusion gap reproduced within 2 points; \method{} with three seeds on RGBT-Tiny & image detector and tables \\
2, uploaded video & identical detections to offline evaluation; TrackEval end to end & video analyser and app \\
3, live streams & 30\,FPS for ten minutes without drift; recovery from unplugging; Orin at 15\,FPS & live system \\
4, multi-camera & counts match manual counts; privacy review signed & anonymous analytics \\
\end{tabularx}
\end{table}
